\subsubsection{A fourth-power bound for scalar bit work}
\label{sec:fixedprecision}

The scalar probabilities can be refined without carrying a growing
common rational denominator.  We give the details because this reduces
the bit overhead in the network theorem.

\begin{lemma}[Fixed-precision coefficient evaluation]
\label{lem:coefficient-four}
Let \(0<\rho\le2\) be a dyadic rational whose description has at
most \(B\) bits.  For integers \(K\ge0\) and \(m\ge1\), the
coefficient \(d_K\) of
\[
  \sech^2(\rho\sqrt{1-z})=\sum_{j\ge0}d_jz^j
\]
has a certified dyadic enclosure of width \(2^{-m}\), computable
in \(O((B+K+m+1)^4)\) bit operations.  The constant is absolute,
and schoolbook multiplication and division suffice.
\end{lemma}

\paragraph{Proof idea.}
Compute the coefficients of the reciprocal denominator first, then invert
its triangular coefficient system. Explicit error bounds allow rounding
after every operation, so the rational denominators never grow with the
number of operations.

\begin{proof}
It suffices to construct an approximation with error \(2^{-m-2}\).
In the following calculation replace the requested precision \(m\)
by \(m+2\) for this purpose.  Write
\[
 A(z)=\cosh^2(\rho\sqrt{1-z})=\sum_{j\ge0}a_jz^j.
\]
The coefficients have alternating signs before the constant term,
so the power series for cosh gives
\begin{equation}\label{eq:coefficient-norm}
 a_0\ge1,\qquad
 \sum_{j\ge0}|a_j|=\cosh^2(\rho\sqrt2)<1024.
\end{equation}
Indeed, \(\rho\sqrt2<3\), \(e<3\), and
\(\cosh^2 3<e^6<729\).  To bound the reciprocal, write
\(\sqrt{1-z}=u+iv\).  On \(|z|\le1\),
\[
 |z|^2=(u^2+v^2-1)^2+4v^2,
 \qquad |v|\le1/2.
\]
Consequently
\[
 |A(z)|=\sinh^2(\rho u)+\cos^2(\rho v)
 \ge\cos^2 1\ge1/4,
\]
and hence
\begin{equation}\label{eq:reciprocal-coefficient-bound}
 |d_j|\le4\qquad(j\ge0)
\end{equation}
by Cauchy's coefficient estimate.

Define integer precision parameters
\begin{align*}
 T&=m+12(K+1)+\lceil\log_2(K+1)\rceil+20,\\
 M&=T+32,\\
 P&=T+8M+10\lceil\log_2(M+1)\rceil+64.
\end{align*}
All operations below round to \(P\) binary fractional places, with
absolute error at most \(2^{-P}\) per rounding.  Their integer
parts have \(O(M)\) bits.  The input \(\rho\) is retained exactly.

For \(0\le j\le M\), compute
\[
 c_j=\frac{(2\rho)^{2j}}{(2j)!},\qquad
 c_0=1,\qquad
 c_{j+1}=c_j\frac{4\rho^2}{(2j+1)(2j+2)}.
\]
Use these values to form the coefficients through degree \(K\) of
\[
 A_M(z)=\frac12\left(1+\sum_{j=0}^M c_j(1-z)^j\right).
\]
The binomial coefficients are exact integers, obtained by Pascal's
recurrence; each has at most \(M+1\) bits.  Terms of degree above
\(K\) may be discarded, but this is unnecessary for the stated
complexity.

The multipliers in the recurrence for \(c_j\) are at most 16.
Its accumulated rounding error is at most
\(16^{M+1}2^{-P}\) at every index.  Expanding the binomials adds
at most a factor \((M+1)2^M\).  The remaining additions and
roundings add at most \((M+1)^2 2^{-P}\) to a coefficient error.
The specified value of \(P\) bounds their sum by
\(2^{-T-4}\).

The truncation error on \(|z|\le1\) is at most
\[
 E_M=\frac12\sum_{j>M}\frac{32^j}{(2j)!}
 \le2^{32-4M}\le2^{-T-4}.
\]
For the middle inequality, \(M\ge16\); successive terms have
ratio below \(1/16\), and
\(34!\ge2^{118}\) follows from
\(16!\ge2^{44}\), \(\prod_{j=17}^{32}j\ge2^{64}\), and
\(33\cdot34\ge2^{10}\).
Cauchy's estimate applies to each coefficient error.
Thus the computed coefficients satisfy
\begin{equation}\label{eq:coefficient-input-error}
 |\widehat a_j-a_j|\le\Delta:=2^{-T},\qquad 0\le j\le K,
 \qquad \widehat a_0\ge1/2.
\end{equation}

Now use the reciprocal recurrence with the same precision:
\[
 \widehat d_0=\operatorname{round}_P(1/\widehat a_0),\qquad
 \widehat d_k=\operatorname{round}_P\left(
 -\frac{\sum_{j=1}^k\widehat a_j\widehat d_{k-j}}
        {\widehat a_0}\right).
\]
Each product is rounded before summation.  Put
\(E_k=\max_{0\le j\le k}|\widehat d_j-d_j|\).
Provided \(E_{k-1}\le1\), equations
\eqref{eq:coefficient-norm}--\eqref{eq:coefficient-input-error}
give
\[
 E_0\le3\Delta,\qquad
 E_k\le2048E_{k-1}+(12K+8193)\Delta\quad(1\le k\le K).
\]
To see the second estimate, the error in the numerator is at most
\(1024E_{k-1}+5k\Delta+k2^{-P}\).  Division by
\(\widehat a_0\) multiplies this by at most two.  Changing
\(1/a_0\) to \(1/\widehat a_0\) contributes at most
\(8192\Delta\), since the exact numerator is bounded by 4096.
The final rounding contributes at most \(2^{-P}\).

Iterating the affine recurrence gives
\[
 E_k\le\left(3+\frac{12K+8193}{2047}\right)2^{11k}\Delta
 \le2^{11k+14}(K+1)\Delta\quad(0\le k\le K).
\]
The choice of \(T\) makes this less than one at every index and
gives \(E_K<2^{-m}\), validating the assumed bound throughout.
The known error bounds supply a certified interval about
\(\widehat d_K\).

There are \(O(M^2+K^2)\) arithmetic operations.  Every operand
has \(O(B+P+M)=O(B+K+m+1)\) bits, including exact binomial
coefficients, products with the input dyadic rational, and rounded
quotients.  Schoolbook arithmetic therefore costs
\(O((B+K+m+1)^4)\) bit operations in total.
\end{proof}

The elementary exponential series on bounded real arguments can
be evaluated by the same fixed-precision method in cubic bit work:
use linearly many terms, linearly many guard bits, and round after
each multiplication or division.  Their total amplification is
bounded by a constant on the argument range used here.  A fourth-power
bound therefore covers tanh, the bias probabilities, and their
rational combinations.  Nonzero denominators require \(O(\bits)\)
guard bits as established above.  The exact integers defining
\(C_K\), and the factor \((5/4)^K\) in the arity acceptance
probability, require \(O(K)\) additional bits.

It follows from Lemma~\ref{lem:coefficient-four} that the accepted
arity can be drawn with \(O(\bits^4)\) expected scalar bit work.
Both the proposal arity and the precision stage have geometric
tails, uniformly in the row and the context.  The mean number of
proposals is 60.  All other scalar work at a reached call obeys the
same bound.  Lemma~\ref{lem:stopping} then gives
\begin{equation}\label{eq:fourth-power-overhead}
 \text{expected online bit work}
 =O\!\left(\bits^4\E N_{\rm calls}\right).
\end{equation}
The preprocessing and storage bounds are unchanged.
